% GENERATED FILE. Edit canonical lessons, then run npm run generate:books.
% canonical-source-hash: fnv1a64:66f8bfb52ba8effd
% canonical-lessons: MR-C411-saphar, MR-C411-rapet, MR-C411-pherphatka, MR-C411-hindne, MR-C411-rengalne, MR-R411-review-404-411

\chapter{Trips and Outings}
\label{ch:mr-trips-and-outings}
\hlchaptermodality{411}

\begin{chapteropening}
\textbf{By the end of this chapter:} \emph{I can talk about trips, rides and strolls, and say I roamed or lingered.}

Complete the last lesson of chapter 411: I can talk about trips, rides and strolls, and say I roamed or lingered.
\end{chapteropening}

\section[\texorpdfstring{\mr{सफर} (saphar)}{सफर (saphar)}]{\mr{सफर} (saphar) — a trip; a tour}

\label{lesson:MR-C411-saphar}



\begin{quote}
Before the new one: say the Marathi for a peninsula, then the Marathi for a beachfront.
\end{quote}

\subsection*{You'll want to know: \mr{सफर}}
\textbf{\mr{सफर}} — \emph{saphar} — \textquotedblleft{}a trip; a tour\textquotedblright{}.

\begin{quote}
\mr{आम्ही} \mr{होडीने} \mr{बेटाची} \mr{सफर} \mr{केली}.
\end{quote}

\begin{quote}
\emph{āmhī hoḍīne beṭācī saphar kelī.}
\end{quote}

\begin{quote}
We took a trip to the island by boat.
\end{quote}

Feminine.

\mr{सफर} \mr{करणे}, to make a trip; it sounds light and happy.

\begin{cousinweb}
\textbf{\mr{सफर}} (\emph{saphar}) is a loan through Persian, usually traced to Arabic \emph{safar}, \textquotedblleft{}a journey\textquotedblright{}.
\end{cousinweb}

\begin{grammarlens}[title={What you've built}]
One more word for getting around.
\end{grammarlens}

\subsection*{Your turn}
\begin{itemize}
\raggedright
  \item \emph{Say it:} \emph{saphar}
  \item \emph{Say it:} \emph{saphar}, once more
  \item \emph{Say it:} \emph{āmhī hoḍīne beṭācī saphar kelī.} — We took a trip to the island by boat.
  \item \emph{Recall:} say the Marathi for a peninsula, then the Marathi for a beachfront, then say \emph{saphar} again
\end{itemize}

\begin{tcolorbox}[breakable,colback=teal!4,colframe=teal!35!black,title={Before you move on}]
What does \mr{सफर} mean? (\textquotedblleft{}A trip; a tour\textquotedblright{}.) Say it once more.
\end{tcolorbox}
\section[\texorpdfstring{\mr{रपेट} (rapeṭ)}{रपेट (rapeṭ)}]{\mr{रपेट} (rapeṭ) — a ride; an outing}

\label{lesson:MR-C411-rapet}



\begin{quote}
Before the new one: say the Marathi for a beachfront, then the Marathi for a trip.
\end{quote}

\subsection*{You'll want to know: \mr{रपेट}}
\textbf{\mr{रपेट}} — \emph{rapeṭ} — \textquotedblleft{}a ride; an outing\textquotedblright{}.

\begin{quote}
\mr{सकाळी} \mr{आम्ही} \mr{सायकलवरून} \mr{रपेट} \mr{मारली}.
\end{quote}

\begin{quote}
\emph{sakāḷī āmhī sāykalvarūn rapeṭ mārlī.}
\end{quote}

\begin{quote}
In the morning we went for a ride on our bicycles.
\end{quote}

Feminine.

\mr{रपेट} \mr{मारणे}, to go for a quick ride or round.

\begin{cousinweb}
\textbf{\mr{रपेट}} (\emph{rapeṭ}) is a Marathi word for a short ride or round trip; its origin is not settled.
\end{cousinweb}

\begin{grammarlens}[title={What you've built}]
One more word for getting around.
\end{grammarlens}

\subsection*{Your turn}
\begin{itemize}
\raggedright
  \item \emph{Say it:} \emph{rapeṭ}
  \item \emph{Say it:} \emph{rapeṭ}, once more
  \item \emph{Say it:} \emph{sakāḷī āmhī sāykalvarūn rapeṭ mārlī.} — In the morning we went for a ride on our bicycles.
  \item \emph{Recall:} say the Marathi for a beachfront, then the Marathi for a trip, then say \emph{rapeṭ} again
\end{itemize}

\begin{tcolorbox}[breakable,colback=teal!4,colframe=teal!35!black,title={Before you move on}]
What does \mr{रपेट} mean? (\textquotedblleft{}A ride; an outing\textquotedblright{}.) Say it once more.
\end{tcolorbox}
\section[\texorpdfstring{\mr{फेरफटका} (pherphaṭkā)}{फेरफटका (pherphaṭkā)}]{\mr{फेरफटका} (pherphaṭkā) — a stroll}

\label{lesson:MR-C411-pherphatka}



\begin{quote}
Before the new one: say the Marathi for a trip, then the Marathi for a ride.
\end{quote}

\subsection*{You'll want to know: \mr{फेरफटका}}
\textbf{\mr{फेरफटका}} — \emph{pherphaṭkā} — \textquotedblleft{}a stroll\textquotedblright{}.

\begin{quote}
\mr{जेवणानंतर} \mr{आम्ही} \mr{बागेत} \mr{फेरफटका} \mr{मारला}.
\end{quote}

\begin{quote}
\emph{jevṇānantar āmhī bāget pherphaṭkā mārlā.}
\end{quote}

\begin{quote}
After dinner we took a stroll in the garden.
\end{quote}

Masculine.

\mr{फेरफटका} \mr{मारणे}, to take a stroll.

\begin{cousinweb}
\textbf{\mr{फेरफटका}} (\emph{pherphaṭkā}) is an echo compound built on \textbf{\mr{फेर}}, \textquotedblleft{}a turn, a round\textquotedblright{}, from the family of \emph{*phirati}, \textquotedblleft{}turns\textquotedblright{} (CDIAL 9078), with \emph{phaṭkā} added for rhythm.
\end{cousinweb}

\begin{grammarlens}[title={What you've built}]
One more word for getting around.
\end{grammarlens}

\subsection*{Your turn}
\begin{itemize}
\raggedright
  \item \emph{Say it:} \emph{pherphaṭkā}
  \item \emph{Say it:} \emph{pherphaṭkā}, once more
  \item \emph{Say it:} \emph{jevṇānantar āmhī bāget pherphaṭkā mārlā.} — After dinner we took a stroll in the garden.
  \item \emph{Recall:} say the Marathi for a trip, then the Marathi for a ride, then say \emph{pherphaṭkā} again
\end{itemize}

\begin{tcolorbox}[breakable,colback=teal!4,colframe=teal!35!black,title={Before you move on}]
What does \mr{फेरफटका} mean? (\textquotedblleft{}A stroll\textquotedblright{}.) Say it once more.
\end{tcolorbox}
\section[\texorpdfstring{\mr{हिंडणे} (hiṇḍṇe)}{हिंडणे (hiṇḍṇe)}]{\mr{हिंडणे} (hiṇḍṇe) — to roam}

\label{lesson:MR-C411-hindne}



\begin{quote}
Before the new one: say the Marathi for a ride, then the Marathi for a stroll.
\end{quote}

\subsection*{You'll want to know: \mr{हिंडणे}}
\textbf{\mr{हिंडणे}} — \emph{hiṇḍṇe} — \textquotedblleft{}to roam\textquotedblright{}.

\begin{quote}
\mr{सुट्टीत} \mr{आम्ही} \mr{गावभर} \mr{हिंडलो}.
\end{quote}

\begin{quote}
\emph{suṭṭīt āmhī gāvbhar hiṇḍlo.}
\end{quote}

\begin{quote}
In the holidays we roamed all over the village.
\end{quote}

Past (he …): \mr{तो} \mr{हिंडला}, \emph{to hiṇḍlā}.

It needs no goal: \mr{इकडेतिकडे} \mr{हिंडणे}, to roam here and there.

\begin{cousinweb}
\textbf{\mr{हिंडणे}} (\emph{hiṇḍṇe}) is inherited: CDIAL lists Marathi \emph{hĩḍṇẽ}, \textquotedblleft{}to move, walk\textquotedblright{}, under Sanskrit \emph{hiṇḍate}, \textquotedblleft{}roams\textquotedblright{} (CDIAL 14089).
\end{cousinweb}

\begin{grammarlens}[title={What you've built}]
One more word for getting around.
\end{grammarlens}

\subsection*{Your turn}
\begin{itemize}
\raggedright
  \item \emph{Say it:} \emph{hiṇḍṇe}
  \item \emph{Say it:} \emph{hiṇḍṇe}, once more
  \item \emph{Say it:} \emph{suṭṭīt āmhī gāvbhar hiṇḍlo.} — In the holidays we roamed all over the village.
  \item \emph{Recall:} say the Marathi for a ride, then the Marathi for a stroll, then say \emph{hiṇḍṇe} again
\end{itemize}

\begin{tcolorbox}[breakable,colback=teal!4,colframe=teal!35!black,title={Before you move on}]
What does \mr{हिंडणे} mean? (\textquotedblleft{}To roam\textquotedblright{}.) Say it once more.
\end{tcolorbox}
\section[\texorpdfstring{\mr{रेंगाळणे} (reṅgāḷṇe)}{रेंगाळणे (reṅgāḷṇe)}]{\mr{रेंगाळणे} (reṅgāḷṇe) — to linger}

\label{lesson:MR-C411-rengalne}



\begin{quote}
Before the new one: say the Marathi for a stroll, then the Marathi for to roam.
\end{quote}

\subsection*{You'll want to know: \mr{रेंगाळणे}}
\textbf{\mr{रेंगाळणे}} — \emph{reṅgāḷṇe} — \textquotedblleft{}to linger\textquotedblright{}.

\begin{quote}
\mr{मुलं} \mr{बागेत} \mr{उशिरापर्यंत} \mr{रेंगाळली}.
\end{quote}

\begin{quote}
\emph{mulaṁ bāget uśirāparyant reṅgāḷlī.}
\end{quote}

\begin{quote}
The children lingered in the park until late.
\end{quote}

Past (he …): \mr{तो} \mr{रेंगाळला}, \emph{to reṅgāḷlā}.

It hints that you should already be gone.

\begin{cousinweb}
\textbf{\mr{रेंगाळणे}} (\emph{reṅgāḷṇe}) is a Marathi verb for staying on longer than you should; its origin is not settled.
\end{cousinweb}

\begin{grammarlens}[title={What you've built}]
One more word for getting around.
\end{grammarlens}

\subsection*{Your turn}
\begin{itemize}
\raggedright
  \item \emph{Say it:} \emph{reṅgāḷṇe}
  \item \emph{Say it:} \emph{reṅgāḷṇe}, once more
  \item \emph{Say it:} \emph{mulaṁ bāget uśirāparyant reṅgāḷlī.} — The children lingered in the park until late.
  \item \emph{Recall:} say the Marathi for a stroll, then the Marathi for to roam, then say \emph{reṅgāḷṇe} again
\end{itemize}

\begin{tcolorbox}[breakable,colback=teal!4,colframe=teal!35!black,title={Before you move on}]
What does \mr{रेंगाळणे} mean? (\textquotedblleft{}To linger\textquotedblright{}.) Say it once more.
\end{tcolorbox}
\section[(dialogue)]{Looking back: from Traffic and Detours to Trips and Outings}

\label{lesson:MR-R411-review-404-411}



\begin{quote}
Say the Marathi for a detour and road traffic.
\end{quote}

\subsection*{Your turn}
Say these aloud:

\begin{itemize}
\raggedright
  \item a detour, road traffic, a traffic signal, to brush against, to crowd thickly
  \item an aeroplane, a flight, a runway, border customs duty, a visa
  \item a pavement, a pedestrian, a signboard, a flyover, to gaze at
  \item a distance, speed, a brake, a signpost, a long trudge
  \item a boatman, an oar, to row, to toss about, an anchor
  \item a sailor, low tide, a sea surge, to grow rough, to slosh about
  \item a coastline, a creek, a lighthouse, a peninsula, a beachfront
  \item a trip, a ride, a stroll, to roam, to linger
\end{itemize}

\begin{tcolorbox}[breakable,colback=teal!4,colframe=teal!35!black,title={Before you move on}]
A detour? (\textbf{\mr{वळसा}}, \emph{vaḷsā}.) Road traffic? (\textbf{\mr{रहदारी}}, \emph{rahdārī}.) A traffic signal? (\textbf{\mr{सिग्नल}}, \emph{signal}.) To brush against? (\textbf{\mr{खेटणे}}, \emph{kheṭṇe}.) To crowd thickly? (\textbf{\mr{दाटणे}}, \emph{dāṭṇe}.) An aeroplane? (\textbf{\mr{विमान}}, \emph{vimān}.) A flight? (\textbf{\mr{उड्डाण}}, \emph{uḍḍāṇ}.) A runway? (\textbf{\mr{धावपट्टी}}, \emph{dhāvpaṭṭī}.) Border customs duty? (\textbf{\mr{सीमाशुल्क}}, \emph{sīmāśulka}.) A visa? (\textbf{\mr{व्हिसा}}, \emph{vhisā}.) A pavement? (\textbf{\mr{पदपथ}}, \emph{padpath}.) A pedestrian? (\textbf{\mr{पादचारी}}, \emph{pādcārī}.) A signboard? (\textbf{\mr{फलक}}, \emph{phalak}.) A flyover? (\textbf{\mr{उड्डाणपूल}}, \emph{uḍḍāṇpūl}.) To gaze at? (\textbf{\mr{न्याहाळणे}}, \emph{nyāhāḷṇe}.) A distance? (\textbf{\mr{अंतर}}, \emph{antar}.) Speed? (\textbf{\mr{वेग}}, \emph{veg}.) A brake? (\textbf{\mr{ब्रेक}}, \emph{brek}.) A signpost? (\textbf{\mr{दिशादर्शक}}, \emph{diśādarśak}.) A long trudge? (\textbf{\mr{पायपीट}}, \emph{pāypīṭ}.) A boatman? (\textbf{\mr{नावाडी}}, \emph{nāvāḍī}.) An oar? (\textbf{\mr{वल्हे}}, \emph{valhe}.) To row? (\textbf{\mr{वल्हवणे}}, \emph{valhavṇe}.) To toss about? (\textbf{\mr{हेलकावणे}}, \emph{helkāvṇe}.) An anchor? (\textbf{\mr{नांगर}}, \emph{nāṅgar}.) A sailor? (\textbf{\mr{खलाशी}}, \emph{khalāśī}.) Low tide? (\textbf{\mr{ओहोटी}}, \emph{ohoṭī}.) A sea surge? (\textbf{\mr{उधाण}}, \emph{udhāṇ}.) To grow rough? (\textbf{\mr{खवळणे}}, \emph{khavaḷṇe}.) To slosh about? (\textbf{\mr{हिंदकळणे}}, \emph{hindkaḷṇe}.) A coastline? (\textbf{\mr{किनारपट्टी}}, \emph{kinārpaṭṭī}.) A creek? (\textbf{\mr{खाडी}}, \emph{khāḍī}.) A lighthouse? (\textbf{\mr{दीपगृह}}, \emph{dīpgṛh}.) A peninsula? (\textbf{\mr{द्वीपकल्प}}, \emph{dvīpkalpa}.) A beachfront? (\textbf{\mr{चौपाटी}}, \emph{caupāṭī}.) A trip? (\textbf{\mr{सफर}}, \emph{saphar}.) A ride? (\textbf{\mr{रपेट}}, \emph{rapeṭ}.) A stroll? (\textbf{\mr{फेरफटका}}, \emph{pherphaṭkā}.) To roam? (\textbf{\mr{हिंडणे}}, \emph{hiṇḍṇe}.) To linger? (\textbf{\mr{रेंगाळणे}}, \emph{reṅgāḷṇe}.)
\end{tcolorbox}
